\section{1970--1990-е: mecha, magical girl, idol, school romance, первые большие базы типажей}

К зрелости 1970--1990-х годов индустрия перестаёт быть россыпью удачных отдельных сериалов и начинает штамповать устойчивые жанровые семьи. Для этой книги важны не сюжеты этих жанров, а то, какую конфигурацию силового поля вокруг персонажа каждый из них закрепляет, — потому что именно эти конфигурации, а не конкретные тайтлы, мы будем находить дальше в главах 2, 3 и 9.

Возьмём первый скачок серьёзно, а не бегло. В "Mazinger Z" (1972) подросток по имени Кодзи Кабуто садится в кабину боевого робота, унаследованного от деда, — и вот в чём фокус этого хода: пилотская кабина внутри головы гигантской машины буквально помещает подростковую психику в центр силы, масштаб которой её пугает и которой она объективно не готова управлять. Всё, что творится внутри — сомнение, паника, минутная слабость, — здесь не остаётся личным делом героя. Оно немедленно и без пересказа становится видно всему миру как масштаб разрушения или спасения. Робот в этой конструкции — не оружие ради оружия. Это буквализированная метафора: что происходит, когда взрослую власть и взрослую ответственность вручают тому, кто ещё психологически подросток, а весь мир вокруг молча ждёт, справится он или нет.

Это первый инструмент, который стоит забрать с собой: сценарист может взять внутреннее состояние — тревогу компетентности, страх подвести, ощущение "мне дали то, к чему я не готов", — и вместо того, чтобы описывать его текстом, буквально построить машину, которая заставляет зрителя пережить это состояние вместе с героем через масштаб последствий. Кабина гигантского робота — это овеществлённая тревога взросления, а не просто дизайнерское решение "давайте сделаем робота покруче".

А дальше, в 1979 году, "Gundam" делает с этим же инструментом нечто более жестокое. Формула не меняется — тот же подросток в той же кабине, — но убирает страховку. В "Mazinger Z" ты почти всегда точно знаешь: это враг, это друг, а поражение — временная неудача, за которой снова придёт победа. В "Gundam" робот твоего героя убивает конкретных, названных, симпатичных зрителю персонажей на другой стороне — а у другой стороны, как выясняется, тоже есть резоны. Тот же самый инструмент — овеществлённая власть в руках подростка — вдруг перестаёт быть метафорой роста и становится метафорой вины. Не "смогу ли я справиться с этой силой", а "что эта сила уже сделала моими руками, и как мне с этим жить дальше". Кокпит остаётся тем же самым местом. Меняется только то, какую именно внутреннюю температуру он теперь измеряет.

Вот первый пункт на карте: один и тот же сценарный инструмент — герой внутри огромной, объективно опасной машины — может обслуживать две совершенно разные внутренние истории, и разница между ними — это разница между power fantasy и war trauma. Для Рин этот инструмент не буквален, но принцип переносится один в один: любая её "сила" — интеллектуальная точность, острота, доступ к системе — точно так же может либо обслуживать историю "я расту и учусь с этим справляться", либо историю "я уже причинила этим вред и должна с этим считаться". Мы вернёмся к этому прямо как к оси в 2.7, когда будем разбирать персонажа как риск.

Magical girl настраивает совершенно другую конфигурацию силового поля. Здесь сила не наследуется снаружи, а рождается изнутри и напрямую зависит от эмоционального состояния — и стоит остановиться и разобрать, что это значит структурно, а не просто перечислить признаки жанра.

Жанр стартует с "Sally the Witch" (1966), но современную формулу — трансформационная последовательность, команда героинь, приватная двойная жизнь — закрепляет "Sailor Moon" (1992) поколением позже. Возьмём саму трансформационную сцену как инструмент, а не как факт из справочника. Каждую серию сюжет буквально останавливается ради фиксированной, неизменной, почти ритуальной последовательности: та же фраза, та же хореография, те же визуальные ритмы — вне зависимости от того, что происходило в серии до этого момента. Почему сериал, построенный на постоянной новизне сюжета, сознательно вставляет в самые драматичные точки нечто неизменное и повторяющееся? Потому что этот ритуал делает то, чего не может сделать сам сюжет: он публично и без права на подделку доказывает, что героиня под костюмом всё ещё честна с собой, всё ещё готова встать и назвать, кто она есть. Трансформация работает не как единоразовое посвящение, а как периодически обновляемое доказательство подлинности, которое приходится заново отыгрывать каждый раз, а не предъявить один раз и больше не возвращаться к вопросу.

Вот второй пункт на карте: сила, привязанная не к тренировке и не к технологии, а к готовности снова и снова, вслух, при свидетелях, называть своё настоящее чувство. Это ранняя, ещё грубая версия принципа, который в главе 3 станет ядром всей dere-теории: там уязвимость и тепло — не слабость персонажа, а именно то, что подтверждает и удерживает его силу.

Idol-аниме даёт третий, самый странный тип силового поля — потому что здесь напряжение возникает не между двумя полюсами вокруг персонажа, а между двумя онтологическими планами, в которых героиня существует одновременно. "Creamy Mami" (1983) и Линн Минмей из "Macross" (1982) — вымышленные певицы внутри сюжета, но их песни выходили как настоящие синглы и звучали по настоящему радио вне сюжета. Один и тот же голос принадлежит одновременно персонажу с внутренней историей (в кадре) и реальному коммерческому артисту (за кадром) — и зазор между этими двумя планами не баг, а именно то, что заставляет феномен работать: слушатель одновременно знает, что это вымысел, и покупает пластинку так, будто это не так.

Вот третий пункт на карте, и он прямо ведёт в 1.6: персонаж может существовать сразу на двух онтологических уровнях — как история и как реальное присутствие — и это не изобретение VTuber-индустрии, а идея, которой уже больше сорока лет. Для RIN это не абстракция: её собственная проектная спецификация уже формализует ровно этот же зазор под именами substrate truth и persona truth — просто сорок лет спустя и по другому поводу.

Школьная романтика не даёт нового полюса силы — она строит сцену и расставляет на ней вспомогательные фигуры, чья работа не в том, чтобы иметь собственную драму, а в том, чтобы держать контур вокруг главной пары. Лучший друг, который знает о чувствах героя раньше него самого и подталкивает к признанию, — функционально тот же вестник, что несёт новости в греческой трагедии. Соперница, которая появляется ровно для того, чтобы заставить героиню назвать вслух, чего она на самом деле хочет, — работает как испытание у порога, знакомое из мифа задолго до школьной формы. Учитель или сосед, который знает больше, чем говорит, и появляется именно в нужный момент, — это хор: не участник действия, а голос, который проговаривает то, что зритель и так подозревал, но персонажи ещё не могут сказать друг другу напрямую.

Эти позиции важно называть по их архетипическому предку, а не просто перечислять как "второстепенных персонажей" — потому что именно так становится видно, зачем они вообще существуют в сюжете: не для того, чтобы у них была своя психология, а для того, чтобы через них сюжет мог сказать то, чего главные герои сказать пока не готовы. В главе 3 мы увидим, как ровно эти же позиции — вестник, испытание, хор — продолжают работать вокруг Рин, только без школьного антуража.

К концу 1990-х на карте уже три разных типа поля: внешняя сила, обслуживающая либо рост, либо вину (mecha); внутренняя сила, требующая постоянного публичного подтверждения (magical girl); персонаж, расколотый между вымыслом и реальным присутствием (idol) — плюс набор вспомогательных архетипических позиций (вестник, испытание, хор), унаследованных из школьной романтики напрямую от античного театра. Это ещё не database character в полном смысле термина (до него ещё десятилетие, см. 1.5), но именно из таких силовых конфигураций, а не из перечня тропов, эта книга будет строить инструментарий дальше.
